\section{Holomorphic 1-Forms and Finite Dimensionality}
\label{sec:holomorphic-forms-finite-dim}

\begin{layman}
Every smooth surface has a notion of \emph{tangent vector} at a
point --- a little arrow giving a direction of motion. Dual to that,
every surface has a notion of \emph{cotangent vector}: a tiny linear
gadget that takes a tangent vector and returns a number. A
\emph{holomorphic 1-form} on a Riemann surface is a global,
complex-linear assignment of such a gadget at every point, varying
holomorphically. In any chart \(z\), a holomorphic 1-form looks like
\(h(z)\,dz\) for some holomorphic function \(h\); under a change of
chart it transforms by the chain rule. The space of all such global
forms is denoted \(H^0(X,\Omega^1_X)\). Its complex dimension is the
\emph{analytic genus} \(g(X)\) --- the most fundamental invariant we
attach to \(X\).

The deep fact of this chapter is that \(g(X)\) is finite. The
argument is a triumph of twentieth-century complex analysis: equip
\(H^0(X,\Omega^1_X)\) with a sup-norm; show that the unit ball is
compact (\emph{Montel's theorem}: a bounded sequence of holomorphic
functions has a uniformly convergent subsequence); apply Riesz's
lemma (a normed space whose unit ball is compact must be
finite-dimensional). The compactness of the surface is what makes
the sup-norm finite in the first place --- on the entire complex
plane, the dimension of holomorphic 1-forms is uncountably infinite.

The chapter develops this in stages. The cotangent bundle gets a
fibre-by-fibre Hermitian norm; that norm assembles into a global
sup-norm by taking the maximum over the (compact) surface; under
uniform-on-compact convergence, sequences of holomorphic 1-forms
behave like sequences of holomorphic functions on a chart, so
Montel's theorem on the chart side propagates to the global form
space; closed unit ball compactness then forces finite
dimensionality through Riesz.

Why bother? Genus is the single most important invariant of a
Riemann surface; it controls everything else. Tying it down
rigorously to a finite-dimensional vector-space dimension is what
lets all the later constructions --- Riemann--Roch, the Jacobian,
Abel--Jacobi, push-pull --- even be \emph{stated}. The next chapter
combines this finite-dim result with sheaf cohomology to prove
Riemann--Roch.
\end{layman}

Let \(\HX\) be the vector space of global holomorphic one-forms. Define the
analytic genus by \(g(X)=\dim_{\C}\HX\), as in
Definition~\ref{def:analytic-genus}.

\subsection{Finite-dimensionality (analytic route)}

The project's preferred route to \(\dim_{\C}\HX<\infty\) is analytic: a
sup-norm makes \(\HX\) a normed space, Montel's theorem makes the closed
unit ball compact, and Riesz's lemma forces finite dimensionality.



\begin{layman}
At each point of the surface, a holomorphic 1-form picks out a tiny
linear gadget --- a local rule that turns a direction into a complex number.
This local gadget lives in a one-dimensional space called the
\emph{cotangent fiber}, and we want to measure how big the gadget is.
The trick: borrow a coordinate chart at the point (a bit of flat graph
paper laid over the curvy surface), use it to declare a unit length on
tangent directions, and then ask, ``what's the largest absolute value the
gadget outputs on a unit-length input?'' That maximum is its norm.
Different charts give slightly different units, but we standardize by
always using the canonical chart that's already attached to the surface's
atlas. Why bother? With a length on every fiber, we can ask whether a
1-form is big or small at any point --- a prerequisite to defining a global
``size'' on the whole space of holomorphic 1-forms in the next step.
\end{layman}



\begin{layman}
We've measured the size of a holomorphic 1-form at each point. To get a
single number summarizing the whole 1-form on the whole surface, just
take the largest pointwise size that ever occurs. This \emph{supremum
norm} is finite because the surface is compact, the size function is
continuous, and a continuous function on a closed-and-bounded space
always attains a maximum (no sneaking off to infinity). The standard
checks pass: the norm is zero exactly for the zero 1-form, scales
correctly under multiplication by a complex number, and obeys the
triangle inequality. So the space of holomorphic 1-forms becomes a
\emph{normed vector space} --- the same structural home as Euclidean space
or any Banach space. Why bother? Once you have a norm, you have notions
of distance, convergence, balls, and compactness. The whole analytic
strategy for proving finite-dimensionality runs on those notions.
\end{layman}


\begin{proof}\leanok\end{proof}

\begin{layman}
Open up a single coordinate chart --- a flat patch of graph paper laid over
a piece of the surface. On that patch, every holomorphic 1-form looks
like an ordinary holomorphic function multiplied by the differential of
the coordinate. The lemma says the size of that ordinary function at
each point is bounded above by a fixed constant --- depending only on the
chart --- times the global supremum norm of the 1-form. In plain terms,
shrinking a 1-form globally shrinks its local coefficient on the chart
by at most a fixed factor; the chart can't blow up the constants behind
your back. Why bother? It's the bridge between the abstract norm on the
space of 1-forms and the concrete behavior of holomorphic functions on
disks in the plane, where Mathlib's existing complex-analysis toolbox
already applies. We need that bridge to invoke Cauchy's estimates and
Montel's theorem in the next step.
\end{layman}

\begin{proofsketch}
Set \(p := \varphi^{-1}(z)\). The chart trivialisation identifies the
fiber \(T^{*}_{p}X\) with \(\C \to_{\C} \C\) isometrically up to a
chart-dependent positive constant \(C_{\varphi}^{-1}\); under this
identification \(\omega_{p}\) corresponds to the linear map
\(\xi\mapsto h_{\varphi}(z)\,\xi\), whose operator norm is exactly
\(|h_{\varphi}(z)|\). Therefore
\(|h_{\varphi}(z)| \le C_{\varphi}\,\|\omega_{p}\|_{p} \le
C_{\varphi}\,\|\omega\|\). The constant \(C_{\varphi}\) is bounded by
the supremum on the (relatively compact) chart image of the
trivialisation comparison factor, which is finite by continuity.
\end{proofsketch}


\begin{proof}\leanok\end{proof}

\begin{layman}
Now the analytic punchline. Take any sequence of holomorphic 1-forms
whose norms are all bounded by 1 --- the closed unit ball. The claim: you
can always extract a subsequence converging uniformly on the whole
surface to another holomorphic 1-form of norm at most 1. This is Montel's
theorem in disguise --- magic special to holomorphic functions, with no
analogue for plain continuous or merely smooth functions. The proof works
chart by chart: in each chart, Cauchy's integral formula bounds
derivatives, which forces the sequence to be \emph{equicontinuous} --- every
form in the sequence wiggles at the same controlled rate. Arzel\`a--Ascoli
then extracts a convergent subsequence, and the chartwise limits glue
into a global holomorphic limit by Weierstrass's convergence theorem
(uniform limits of holomorphic functions are again holomorphic). Why
bother? It says the unit ball is sequentially compact --- a heavy hammer
that, in normed spaces, is the trademark signature of being
finite-dimensional.
\end{layman}

\begin{proofsketch}
\emph{Step 1: finite chart cover.} Compactness of \(X\) supplies a
finite atlas \(\{(U_i,\varphi_i)\}_{i=1}^{N}\) with each
\(\varphi_i(U_i) \supseteq \overline{B(c_i, R_i)}\); choose
\(0 < r_i < R_i\) with \(\bigcup_i \varphi_i^{-1}(B(c_i, r_i)) = X\).

\emph{Step 2: chartwise sup bound.} Lemma~\ref{lem:chart-coefficient-bound}
gives constants \(C_i\) such that the chart coefficient
\(h_{n,i}\) of \(\omega_n\) on \(U_i\) satisfies
\(|h_{n,i}(z)| \le C_i\) for all \(z \in \overline{B(c_i, R_i)}\) and
all \(n\).

\emph{Step 3: Cauchy estimates yield equicontinuity.} For
\(z \in \overline{B(c_i, r_i)}\) the Cauchy integral formula for the
derivative,
\(h_{n,i}'(z) = \frac{1}{2\pi i}\oint_{|w-c_i|=R_i}
\frac{h_{n,i}(w)}{(w-z)^2}\,dw\),
combined with \(|w-z| \ge R_i - r_i\), gives
\(|h_{n,i}'(z)| \le C_i \,R_i / (R_i - r_i)^2\). The family
\(\{h_{n,i}\}\) is therefore uniformly Lipschitz, hence equicontinuous,
on \(\overline{B(c_i, r_i)}\). (Mathlib:
\texttt{DiffContOnCl.circleIntegral\_sub\_inv\_smul} for the
Cauchy formula;
\texttt{circleIntegral.norm\_integral\_le\_of\_norm\_le\_const} for the
\(L^{\infty}\)--to--integral bound.)

\emph{Step 4: Arzel\`a--Ascoli on each chart.} The Arzel\`a--Ascoli theorem
(Mathlib \texttt{ArzelaAscoli.isCompact\_closure\_of\_isClosedEmbedding})
extracts a subsequence converging uniformly on
\(\overline{B(c_i, r_i)}\). Iterating across the \(N\) charts
diagonally produces a single subsequence
\(\omega_{\varphi(n)}\) whose chart restrictions all converge
uniformly.

\emph{Step 5: glue and pass to the limit.} The chartwise uniform
limits agree on overlaps (because they are obtained from a single
diagonal subsequence), so they assemble into a function on \(X\) which
is the uniform limit of \(\omega_{\varphi(n)}\). By Weierstrass's
convergence theorem (uniform limit of holomorphic functions is
holomorphic; cf.\ Forster \S 14 or Griffiths--Harris pp.~7--8) the chart
coefficients of the limit are again holomorphic, so the limit is a
holomorphic 1-form \(\omega_\infty\). The sup-norm bound
\(\|\omega_\infty\|\le 1\) follows by taking the limit in
\(|h_{\varphi(n),i}(z)|\le 1\) (with our normalisation; if not unit
norm, by passing to the limit in the explicit estimate).
\end{proofsketch}


\begin{proof}\leanok\end{proof}

\begin{layman}
The previous lemma extracted convergent subsequences --- that's
\emph{sequential} compactness. This theorem upgrades it to ordinary
\emph{topological} compactness, the textbook ``every open cover has a
finite subcover'' kind. The two flavors agree in any normed space,
because such spaces are metrisable, and in metric spaces sequential and
topological compactness are the same notion. So the closed unit ball of
holomorphic 1-forms --- the set of forms with norm at most 1 --- is compact in
the textbook sense. Why bother? Compactness of a closed ball is one of
those innocent-sounding facts that's actually a structural giveaway: in
a normed space, it \emph{forces} the space to be finite-dimensional.
That's the next step.
\end{layman}

\begin{proofsketch}
A normed space is metrisable, hence first-countable; in a first-countable
space sequential compactness and topological compactness coincide
(Mathlib \texttt{UniformSpace.isCompact\_iff\_isSeqCompact} or
\texttt{IsSeqCompact.isCompact}). Apply
Lemma~\ref{lem:montel-compactness} to obtain sequential compactness of
the closed unit ball, then convert. The closed unit ball is closed in
the metric topology by continuity of the norm, so the conclusion is
exactly \texttt{IsCompact (Metric.closedBall 0 1)}.
\end{proofsketch}


\begin{proof}\leanok\end{proof}

\begin{layman}
Here's a surprising fact called Riesz's lemma: in any normed vector
space, if even a single closed ball of positive radius is compact, then
the whole space must be finite-dimensional. The intuition is that
infinite-dimensional spaces are too roomy to be compact --- you can always
find a sequence of unit-distance-apart points that never converges, and
compactness rules this out, so room runs out. Mathlib already has this
packaged. Apply it to our setting: the closed unit ball of holomorphic
1-forms is compact (just shown), therefore the space of holomorphic
1-forms is finite-dimensional. Why bother? This is the headline
conclusion of the entire chapter. The genus of the surface --- defined as
the dimension of this space --- is now guaranteed to be a finite whole
number. Without it, the whole rest of the project (lattices, Jacobians,
Abel--Jacobi maps) would have nothing to grab onto.
\end{layman}

\begin{proofsketch}
Direct invocation of
\texttt{FiniteDimensional.of\_isCompact\_closedBall} \(\C\) \texttt{one\_pos\;h}
where \(h\) is the conclusion of
Theorem~\ref{thm:hone-unit-ball-compact} for the chosen realisation.
\end{proofsketch}


\begin{proof}\leanok\end{proof}

\begin{layman}
We package the analytic route as a single named input: any normed-space
realization of the global holomorphic 1-forms whose norm matches our
sup norm is automatically finite-dimensional. This bundles the fiber
norm, the sup norm, the local coefficient bound, Montel's compactness
theorem, the upgrade from sequential to topological compactness, and
Riesz's lemma. There's a separate route to the same conclusion via
Riemann--Roch --- listed alongside as its own umbrella --- but the analytic
route is the project's preferred path because its prerequisites are
mostly already in Mathlib. Why bother? Naming the umbrella keeps the
dependency graph clean. Downstream theorems depend on
``finite-dimensionality, by some route,'' rather than on the specific
analytic chain, so a future swap of routes won't disturb anything that
depends on it.
\end{layman}

\begin{proofsketch}
Combine Theorem~\ref{thm:hone-unit-ball-compact}, which gives
\texttt{IsCompact (Metric.closedBall 0 1)} on the realisation \(H\),
with Theorem~\ref{thm:fd-from-riesz}.
\end{proofsketch}

\subsection{Compactness of holomorphic 1-forms: details}








\begin{proof}\leanok\end{proof}

\begin{layman}
This is the headline statement, in its public form: the space of global
holomorphic 1-forms on a compact Riemann surface is finite-dimensional.
Equivalently, the analytic genus --- defined as the complex dimension of
that space --- is a finite natural number. We get it by taking any
project-side normed-space realization of this space, applying the
finite-dimensionality umbrella, and transporting the conclusion back.
Why bother? Genus is the single most important invariant of a Riemann
surface. ``Sphere'' has genus 0, ``donut'' has genus 1, ``two-handled
pretzel'' has genus 2, and so on up. Knowing genus is finite means we
can actually compute with it --- and the finite-dimensionality of the
1-forms is what makes period lattices, Jacobians, and the entire
downstream story even formulable.
\end{layman}

\begin{proofsketch}
Transport \(\HX\) into the canonical sup-norm realisation supplied by
the project-side \texttt{HolomorphicOneFormBanachData}
(Jacobian/HolomorphicForms/CompactRiemannSurface.lean), apply
Theorem~\ref{input:finite-dimensionality} to conclude
finite-dimensionality of the realisation, and transport back along
the realisation isomorphism. (\(\C\)-linear bijections preserve
finite-dimensionality.)
\end{proofsketch}
